\section{Evaluation of Generated Review Comments}
\label{sec:sota-evaluation}

How to score a generated review is contested, and the contest has moved through
three positions.

The first was reference-based text similarity: compare the generated comment with
the comment a human actually wrote, using BLEU or a relative. DeepCRCEval
dismantles this, showing that overlap metrics fail to capture the qualitative
dimensions that make a review useful and that models ranked highly by them are not
ranked highly by humans; the authors replace the metric with a criterion-based
rubric scored by a language model~\cite{hong2024deepcrceval}. The objection is not
merely empirical but structural: a change admits many valid reviews, so agreement
with one of them is a poor proxy for quality.

The second position is reference-free scoring grounded in program analysis.
CRScore constructs pseudo-references from claims and code smells detected by
language models and static analysers, then measures conciseness,
comprehensiveness and relevance against those, reporting the strongest alignment
with human judgement among open-source metrics~\cite{naik2025crscore}. This is
the closest prior work to the evaluation in this thesis in spirit because both
refuse a single human reference and ground scoring in what the code contains.
They differ in instrument: CRScore is one metric pipeline, whereas this thesis
uses a panel of judges scoring an explicit rubric.

The third position is benchmark construction with human validation attached, which
is where the 2025 literature sits. SWR-Bench pairs a thousand verified pull
requests with an evaluation method its authors report as agreeing with human
assessment about ninety percent of the time~\cite{zeng2025swrbench}. CodeFuse-CR-Bench
combines rule-based and model-based scoring over repository-level
instances~\cite{guo2025codefusecrbench}. Khan et al.\ survey code-review
benchmarks and evaluation practice across a decade and catalogue how heterogeneous
the protocols remain~\cite{khan2026crsurvey}. Widyasari et al.\ approach quality
from the human side, finding that a substantial fraction of real review comments
offer a suggestion with no explanation at
all~\cite{widyasari2025explanations}. Davila and Nunes, reviewing modern code review
broadly, note that evaluations of review-supporting approaches are frequently
conducted without human involvement~\cite{davila2021mcrslr}.

Two gaps persist across this body of work. The evaluations that adopt a language
model as the scorer overwhelmingly use a \emph{single} judge, so the reliability of
the instrument is neither measured nor reported; the LLM-as-a-judge literature
reviewed next has established that single-judge scoring is exposed to biases that
panels and agreement statistics are designed to expose. And the rubrics measure
whether a review addresses a topic, not whether what it says about that topic is
true. This distinction becomes decisive when the intervention under test
supplies the review with facts it might state incorrectly.
